\documentclass[12pt,a4paper,oneside]{report}

% =========================================================
% PAQUETES
% =========================================================
\usepackage[utf8]{inputenc}
\usepackage[T1]{fontenc}
\usepackage[spanish,es-nodecimaldot]{babel}
\usepackage{lmodern}
\usepackage{microtype}
\usepackage{geometry}
\geometry{top=2.7cm,bottom=2.7cm,left=3cm,right=2.5cm}
\usepackage{amsmath,amssymb,amsfonts,mathtools,bm}
\usepackage{booktabs}
\usepackage{array}
\usepackage{xcolor}
\usepackage{tikz}
\usepackage{wrapfig} % para wrapfigure
\usepackage{graphicx}
\usepackage{amsthm}
\usetikzlibrary{arrows.meta,positioning,shapes.geometric}
\usepackage{hyperref}
\hypersetup{
    colorlinks=true,
    linkcolor=blue!55!black,
    citecolor=blue!55!black,
    urlcolor=blue!65!black,
    pdftitle={Aplicación del método de Benettin al atractor de Lorenz},
    pdfauthor={Documento tutorial}
}
\urlstyle{same}
\theoremstyle{definition}
\newtheorem{ejemplo}{Ejemplo}[chapter]
\setcounter{secnumdepth}{3}
\setcounter{tocdepth}{2}
\numberwithin{equation}{chapter}

% =========================================================
% COMANDOS
% =========================================================
\newcommand{\R}{\mathbb{R}}
\newcommand{\dd}{\mathrm{d}}
\newcommand{\e}{\mathrm{e}}
\newcommand{\tr}{\operatorname{tr}}
\newcommand{\diag}{\operatorname{diag}}
\newcommand{\norm}[1]{\left\lVert #1\right\rVert}
\newcommand{\abs}[1]{\left|#1\right|}
\newcommand{\trans}{\mathsf{T}}

% =========================================================
% DOCUMENTO
% =========================================================
\begin{document}

\begin{titlepage}
    \centering
    \vspace*{1.8cm}
    {\Huge\bfseries Aplicación del método de Benettin\\[0.25cm] al atractor de Lorenz\par}
    
    \vspace{0.7cm}
    {\Large Cálculo del exponente máximo y del espectro completo\\ de exponentes de Lyapunov\par}
    \vfill
    \begin{minipage}{0.95\textwidth}
        \centering
        Sistema clásico de Lorenz:
        \[
            \sigma=10,\qquad \rho=28,\qquad \beta=\frac{8}{3}.
        \]
        Condiciones iniciales:
        \[
            X_0=(1,1,1)^{\trans},\qquad X'_0=(1.000001,1,1)^{\trans}.
        \]
    \end{minipage}
    \vfill
    
    {\large \today\par}
\end{titlepage}

\tableofcontents
\clearpage

% =========================================================
\chapter*{Introducción}
\addcontentsline{toc}{chapter}{Introducción}
% =========================================================

El sistema de Lorenz constituye uno de los modelos canónicos del caos determinista. Fue introducido como una reducción de las ecuaciones de convección de Rayleigh--Bénard y mostró que un sistema autónomo de pocas variables puede presentar movimiento acotado, aperiódico y extremadamente sensible a sus condiciones iniciales \cite{lorenz1963,kuznetsov2020}.

Los exponentes de Lyapunov cuantifican esta sensibilidad. Describen las tasas exponenciales medias con las que perturbaciones infinitesimales crecen o se contraen a lo largo de una trayectoria. El método de Benettin calcula estas tasas propagando vectores en el espacio tangente y renormalizándolos periódicamente \cite{benettin1980,skokos2010}.

Para recuperar el espectro completo es necesario impedir que todos los vectores tangentes terminen alineados con la dirección de máximo crecimiento. Esto se consigue mediante ortogonalización de Gram--Schmidt o, equivalentemente en aritmética exacta, mediante una factorización QR \cite{semenov2022,gatica2026}.

% =========================================================
\chapter{El sistema de Lorenz}
% =========================================================

Sea
\[
    X(t)= \begin{pmatrix} x(t)\\y(t)\\z(t) \end{pmatrix}\in\R^3, \qquad \dot X=F(X).
\]
El sistema de Lorenz es
\begin{equation}
    \boxed{
    \begin{aligned}
        \dot x &= \sigma(y-x),\\
        \dot y &= x(\rho-z)-y,\\
        \dot z &= xy-\beta z.
    \end{aligned}}
    \label{eq:lorenz}
\end{equation}

A efectos de este análisis utilizaremos los parámetros clásicos que hacen que el sistema se comporte de forma caótica.

\begin{equation}
    \sigma=10,\qquad \rho=28,\qquad \beta=\frac83.
\end{equation}

La no linealidad procede de los productos \(xz\) y \(xy\). Por lo que La divergencia del campo vectorial viene dada por:

\begin{equation}
    \nabla\cdot F = \frac{\partial\dot x}{\partial x} +\frac{\partial\dot y}{\partial y} +\frac{\partial\dot z}{\partial z} =-\sigma-1-\beta=-\frac{41}{3}<0.
\end{equation}

Por la fórmula de Liouville  (generalización matemática ligada al teorema de la divergencia en sistemas dinámicos), establece que la tasa de cambio temporal de un elemento de volumen \(V(t)\) en el espacio de fases es proporcional a la divergencia del sistema multiplicada por el propio volumen \cite{Ehrendorfer2003}. 
El resultado es una ecuación diferencial de la forma:

\begin{equation}
    \frac{dV(t)}{dt} = -(\sigma + 1+ \beta)V(t)    
\end{equation}
Que, al integrarla queda de la forma:
\begin{equation}
    V(t)=V(0)\exp\left[-(\sigma+1+\beta)t\right].
\end{equation}

El flujo es, por tanto, \textbf{disipativo}: el volumen se contrae globalmente aunque existan estiramientos locales en determinadas direcciones.

% =========================================================
\chapter{Puntos de equilibrio y estabilidad mediante autovalores}
% =========================================================

\section{Puntos de Equilibrio}
La resolución algebraica completa de \(F(X_e)=0\) se presenta en el Apéndice~\ref{app:equilibrios}. Los puntos resultantes son
\begin{equation}
    E_0=(0,0,0)
\end{equation}
y, para \(\rho>1\),
\begin{equation}
    E_{\pm}= \left( \pm\sqrt{\beta(\rho-1)}, \pm\sqrt{\beta(\rho-1)}, \rho-1 \right).
\end{equation}
Con los parámetros clásicos,
\begin{equation}
    \boxed{ E_{\pm}=(\pm6\sqrt2,\pm6\sqrt2,27). }
\end{equation}

\section{Estudio de Estabilidad del sistema}
Determinamos la matriz Jacobiana del sistema de Lorenz, obteniéndose: 
\begin{equation}
    \boxed{
    J(x,y,z)=DF(x,y,z)=
    \begin{pmatrix}
        -\sigma & \sigma & 0\\
        \rho-z & -1 & -x\\
        y & x & -\beta
    \end{pmatrix}.}
    \label{eq:jacobiana}
\end{equation}

La estabilidad local se determina resolviendo
\begin{equation}
    \boxed{\det\!\left(J(X_e)-\lambda I\right)=0.}
    \label{eq:problema-espectral}
\end{equation}

Si todos los autovalores tienen \textbf{parte real negativa}, el equilibrio es \textbf{localmente estable}. Una parte \textbf{real positiva} determina una \textbf{dirección inestable}.

\section{Análisis espectral de \(E_0\)}
En el origen,
\begin{equation}
    J(E_0)=
    \begin{pmatrix}
        -10&10&0\\
        28&-1&0\\
        0&0&-\frac83
    \end{pmatrix}.
\end{equation}

El problema característico es
\begin{align}
    0&=\det
    \begin{pmatrix}
        -10-\lambda&10&0\\
        28&-1-\lambda&0\\
        0&0&-\frac83-\lambda
    \end{pmatrix}\\
    &= \left(-\frac83-\lambda\right) \left[(-10-\lambda)(-1-\lambda)-280\right].
\end{align}

Por tanto,
\begin{equation}
    \left(\lambda+\frac83\right) \left(\lambda^2+11\lambda-270\right)=0.
\end{equation}

Los autovalores son
\begin{equation}
    \lambda_{1,2}=\frac{-11\pm\sqrt{1201}}{2}, \qquad \lambda_3=-\frac83,
\end{equation}
es decir,
\begin{equation}
    \boxed{ \lambda_1\simeq11.8277,\quad \lambda_2\simeq-22.8277,\quad \lambda_3\simeq-2.6667.}
\end{equation}

El autovalor positivo genera una variedad inestable \textbf{unidimensional}; los dos autovalores negativos generan una variedad estable \textbf{bidimensional}. El origen es, en consecuencia, un punto silla hiperbólico.

\section{Análisis espectral de \(E_{\pm}\)}
Sea
\[
    a=\sqrt{\beta(\rho-1)}.
\]
En \(E_+\),
\begin{equation}
    J(E_+)-\lambda I=
    \begin{pmatrix}
        -\sigma-\lambda&\sigma&0\\
        1&-1-\lambda&-a\\
        a&a&-\beta-\lambda
    \end{pmatrix}.
\end{equation}

Al desarrollar el determinante,
\begin{align}
    0={}&(-\sigma-\lambda) \left[(-1-\lambda)(-\beta-\lambda)+a^2\right]\nonumber\\
    &-\sigma\left[-\beta-\lambda+a^2\right].
\end{align}

Usando \(a^2=\beta(\rho-1)\), se obtiene
\begin{equation}
    \boxed{ \lambda^3+(\sigma+\beta+1)\lambda^2 +\beta(\sigma+\rho)\lambda +2\sigma\beta(\rho-1)=0.}
    \label{eq:polinomio-Epm}
\end{equation}

Debido a que los puntos de Equilibrio son simétricos, entonces tanto \(E_-\) como \(E_+\) poseen el mismo polinomio característico. Así:

\begin{equation}
    \lambda^3+\frac{41}{3}\lambda^2 +\frac{304}{3}\lambda+1440=0.
\end{equation}

Sus raíces aproximadas son:
\begin{equation}
    \boxed{ \lambda_1=-13.8546,\qquad \lambda_{2,3}=0.0940\pm10.1945\,\mathrm{i}.}
\end{equation}

El autovalor real \textbf{negativo representa contracción en una dirección}. El par complejo tiene parte real positiva y produce rotación acompañada de expansión en un plano local. Así, \(E_\pm\) son puntos--silla \textbf{inestables}. Estos autovalores calculados no deben confundirse con los exponentes de Lyapunov del atractor que calcularemos más adelante. Los primeros describen el flujo cerca de un equilibrio; los segundos son promedios asintóticos calculados a lo largo de una trayectoria variable.

% =========================================================
\chapter{Series de Taylor y linealización}
% =========================================================

Considéremos una trayectoria nominal \(X(t)\) y una trayectoria próxima:
\[
    X'(t)=X(t)+\delta X(t).
\]

La expansión de la serie de Taylor de primer orden quedaría expresada:
\begin{equation}
    F(X+\delta X) = F(X)+DF(X)[\delta X]+\mathcal{O}(\|\delta X\|^2).
\end{equation}
Donde \(DF(X)[\cdot]\) es la diferencial (derivada direccional) en \(X\), cuya representación matricial en la base canónica es la matriz jacobiana del sistema \(J(X)\), es decir, \(DF(X)[\delta X]=J(X)\,\delta X\) y partiendo de las dos trayectorias \(X(t)\) e \(X'(t)\) con \(\dot{X}=F(X)\) y \(\dot{X'}=F(X')\), al derivar la perturbación se obtiene:

\begin{equation}
    \dot{\delta X} = \dot{X'}-\dot{X} = F(X')-F(X) = F(X+\delta X)-F(X) = J(X(t))\,\delta X+\mathcal{O}(\|\delta X\|^2).
\end{equation}

En el límite infinitesimal,
\begin{equation}
    \boxed{\delta\dot X=J(X(t))\delta X.}
    \label{eq:variacional}
\end{equation}

El modelo de Lorenz queda expresado,
\begin{equation}
    \begin{aligned}
        \delta\dot x&=-\sigma\delta x+\sigma\delta y,\\
        \delta\dot y&=(\rho-z(t))\delta x-\delta y-x(t)\delta z,\\
        \delta\dot z&=y(t)\delta x+x(t)\delta y-\beta\delta z.
    \end{aligned}
\end{equation}

La Jacobiana cambia con el tiempo porque se evalúa sobre \(X(t)\). Por ello, sus auto-valores instantáneos \textbf{no son, en general, exponentes de Lyapunov}.

% =========================================================
\chapter{Variacional y dinámica tangente}
% =========================================================

Partimos que la matriz fundamental \(\Phi(t,t_0)\) satisface:
\begin{equation}
    \dot\Phi(t,t_0)=J(X(t))\Phi(t,t_0), \qquad \Phi(t_0,t_0)=I.
    \label{eq:fundamental}
\end{equation}

\begin{wrapfigure}{r}{0.55\textwidth}
  \centering
  \vspace{-6pt}
  \includegraphics[width=0.95\linewidth]{img/delta_tangente.png}
  \caption{Separación infinitesimal $\delta(t)$ entre trayectoria de referencia $\mathbf{x}(t)$ y trayectoria perturbada $\mathbf{x}(t)+\delta(t)$. El vector tangente evoluciona linealmente según \eqref{eq:fundamental}.}
  \label{fig:delta-tangente}
  \vspace{-8pt}
\end{wrapfigure}

La perturbación evoluciona según
\begin{equation}
    \delta X(t)=\Phi(t,t_0)\delta X(t_0).
    \label{eq:propagador}
\end{equation}

Geométricamente, \(\Phi\) transforma una \textbf{esfera infinitesimal} de condiciones iniciales en un elipsoide\footnote{En la sección 9.3 de \textit{Nonlinear Dynamics and Chaos} (2018), p.~29, Strogatz introduce la evolución de una pequeña esfera de condiciones iniciales en el espacio de fases, la cual se deforma hasta convertirse en un elipsoide a medida que el flujo avanza.}. Sus ejes representan direcciones tangentes de estiramiento o contracción. La dinámica tangente conserva únicamente la parte lineal de la separación y es válida mientras la perturbación permanezca suficientemente pequeña \cite{benettin1980}. En la Fig.~\ref{fig:delta-tangente} se ilustra $\delta(t)$ como conector entre ambas órbitas, tangente al flujo.

% =========================================================
\chapter{Exponentes de Lyapunov}
% =========================================================

Sea \(v\neq0\), el exponente asociado a la dirección inicial \(v\) es
\begin{equation}
    \lambda(X_0,v)= \limsup_{t\to\infty} \frac{1}{t} \ln \frac{\norm{\Phi(t,0)v}}{\norm{v}}.
\end{equation}

Formalmente,
\begin{equation}
    \norm{\delta X(t)} \sim \norm{\delta X(0)}\e^{\lambda t}.
\end{equation}

Un exponente positivo indica crecimiento exponencial medio; uno negativo, contracción; y uno nulo, una dirección neutra. En un flujo autónomo no estacionario aparece normalmente un exponente cero asociado a la dirección del propio flujo.

Para el Lorenz clásico, un espectro de referencia encontrado en la literatura es:
\begin{equation}
    (\lambda_1,\lambda_2,\lambda_3) \simeq(0.906,0,-14.572)
    \label{eq:espectro-ref}
\end{equation}
\cite{sprott2005}. 

Además,
\begin{equation}
    \sum_{i=1}^{3}\lambda_i = \lim_{T\to\infty}\frac1T \int_0^T\tr J(X(t))\,\dd t =-\sigma-1-\beta=-\frac{41}{3}.
    \label{eq:suma}
\end{equation}

Esta identidad constituye una comprobación de consistencia muy útil.



% =========================================================
\chapter{Método de Benettin}
% =========================================================

La objetivo de principal método de Benettin es medir la divergencia
exponencial media entre dos trayectorias infinitesimalmente próximas.
En un sistema \textbf{caótico}, una separación inicial de norma $\varepsilon$
crece en promedio como $\sim\varepsilon e^{\lambda_1 t}$, donde
$\lambda_1$ es el exponente de Lyapunov máximo. Integrar directamente
durante tiempos largos no funciona: la perturbación sobresale del régimen
lineal (saturación no lineal) o produce desbordamiento numérico. La
solución consiste en propagar por intervalos cortos $\Delta T$ y, al final de
cada uno, conservar únicamente la dirección de la perturbación
restaurando su norma a $\varepsilon$. El exponente se recupera entonces
como el promedio temporal de los logaritmos de los factores de
crecimiento $g_k$ de cada intervalo \cite{benettin1980}.

Sean
\[
    X_0=(1,1,1)^{\trans},\qquad X'_0=(1.000001,1,1)^{\trans}.
\]
Las condiciones iniciales del sistema de Lorenz \textbf{nominal y perturbado} respectivamente; la separación inicial es:
\begin{equation}
    \delta X_0=X'_0-X_0 =10^{-6}(1,0,0)^{\trans}, \qquad \varepsilon=10^{-6}.
\end{equation}

Definimos además intervalos \(t_k=k\Delta T\). Durante cada intervalo se propagan simultáneamente la trayectoria nominal y la trayectoria perturbada (ecuación variacional). Si \(\widetilde{\delta X}_{k+1}\) es la perturbación antes de renormalizar, el factor de crecimiento es:

\begin{equation}
    g_k=\frac{\norm{\widetilde{\delta X}_{k+1}}}{\varepsilon}.
\end{equation}

A continuación se conserva la dirección y se restaura la norma:
\begin{equation}
    \delta X_{k+1} = \varepsilon \frac{\widetilde{\delta X}_{k+1}} {\norm{\widetilde{\delta X}_{k+1}}}.
\end{equation}

La Figura~\ref{fig:benettin_gram_schmidt} ilustra el significado
geométrico de este procedimiento: la dinámica tangente deforma los
vectores y los alinea progresivamente con la dirección dominante,
mientras que la \textbf{reortonormalización} restaura una base ortonormal
conservando únicamente la información de estiramiento. La figura
muestra ya el caso de tres vectores con factorización
$\widetilde W = QR$, cuya formalización para el espectro completo se
desarrolla en la sección siguiente; para el exponente máximo basta
con el primer vector.

\shorthandoff{>}
\begin{figure}[htb]
    \centering
    \resizebox{\textwidth}{!}{%
    \begin{tikzpicture}[
        >=Latex,
        vector/.style={->, line width=1.2pt},
        labelstyle/.style={font=\small},
        every node/.style={font=\small}
    ]
        % =====================================================
        % PANEL 1: BASE ORTONORMAL INICIAL
        % =====================================================
        %\draw[gray!60] (-0.8,-0.8) rectangle (3.2,4.2);
        \node[font=\bfseries] at (1.2,3.7) {Base ortonormal inicial};

        \coordinate (O1) at (1.0,0.8);
        \draw[fill=black] (O1) circle (1.2pt);
        \draw[vector,red] (O1) -- ++(0.2,1.8) node[above] {$q_1$};
        \draw[vector,blue] (O1) -- ++(1.7,1.2) node[right] {$q_2$};
        \draw[vector,green!60!black] (O1) -- ++(1.3,0.2) node[right] {$q_3$};
        \node[align=center] at (1.2,-0.25) {Vectores ortogonales\\normalizados};

        % =====================================================
        % FLECHA PROPAGACION
        % =====================================================
        \draw[->,thick] (3.6,1.8) -- (5.4,1.8);
        \node[align=center] at (4.5,2.5) {Propagación\\durante $\Delta T$};
        \node at (4.5,1.1) {$\dot W = J(X(t))W$};

        % =====================================================
        % PANEL 2: VECTORES DEFORMADOS
        % =====================================================
        %\draw[gray!60] (5.8,-0.8) rectangle (10.0,4.2);
        \node[font=\bfseries] at (7.9,3.7) {Vectores propagados};

        \coordinate (O2) at (6.8,0.8);
        \draw[fill=black] (O2) circle (1.2pt);
        \draw[vector,red] (O2) -- ++(0.6,2.0) node[above] {$\widetilde w_1$};
        \draw[vector,blue] (O2) -- ++(0.8,1.8) node[right] {$\widetilde w_2$};
        \draw[vector,green!60!black] (O2) -- ++(1.0,1.5) node[right] {$\widetilde w_3$};
        \draw[dashed,gray] (O2) -- ++(0.9,1.9);
        \node[align=center] at (7.9,-0.25) {Alineamiento progresivo\\hacia la dirección dominante};

        % =====================================================
        % FLECHA QR (Ajustada para que el texto no toque la línea)
        % =====================================================
        \draw[->,thick] (10.4,1.8) -- (12.2,1.8);
        \node[align=center] at (11.3,2.65) {Ortogonalización\\de Gram--Schmidt};
        \node at (11.3,1.1) {$\widetilde W = QR$};

        % =====================================================
        % PANEL 3: BASE REORTOGONALIZADA
        % =====================================================
        %\draw[gray!60] (12.6,-0.8) rectangle (16.8,4.2);
        \node[font=\bfseries] at (14.7,3.7) {Base reortogonalizada};

        \coordinate (O3) at (13.6,0.8);
        \draw[fill=black] (O3) circle (1.2pt);
        \draw[vector,red] (O3) -- ++(0.3,1.9) node[above] {$q_1^{(k+1)}$};
        \draw[vector,blue] (O3) -- ++(1.9,1.1) node[right] {$q_2^{(k+1)}$};
        \draw[vector,green!60!black] (O3) -- ++(1.4,0.3) node[right] {$q_3^{(k+1)}$};
        \node[align=center] at (14.7,-0.25) {Nueva base ortonormal\\para el siguiente ciclo};

        % =====================================================
        % ACUMULACION DE FACTORES
        % =====================================================
        \draw[rounded corners,fill=blue!5] (6.0,-2.6) rectangle (10.4,-1.2);
        \node[align=center] at (8.2,-1.9) {Factores de estiramiento\\extraídos de $R$};

        \draw[rounded corners,fill=green!5] (11.4,-2.6) rectangle (16.6,-1.2);
        \node[align=center] at (14.0,-1.9) {$\ln|R_{11}|,\;\ln|R_{22}|,\;\ln|R_{33}|$};

        \draw[->,thick] (10.4,-1.9) -- (11.4,-1.9);

        % =====================================================
        % LEYENDA
        % =====================================================
        \node[align=justify,text width=14cm,font=\small] at (8.0,-4.0) {
            \textbf{Interpretación geométrica:} la propagación mediante la dinámica tangente hace que los vectores se alineen con la dirección asociada al mayor exponente de Lyapunov. La ortogonalización de Gram--Schmidt elimina las componentes ya capturadas por los vectores dominantes y reconstruye una base ortonormal, permitiendo estimar todo el espectro de Lyapunov.
        };
    \end{tikzpicture}
    }
    \caption{
        Representación geométrica del método de Benettin con ortogonalización de Gram--Schmidt. Los vectores tangentes se propagan durante un intervalo $\Delta T$ mediante la ecuación variacional, se deforman y tienden a alinearse con la dirección dominante. La reortogonalización recupera una base ortonormal y los elementos diagonales de la matriz $R$ proporcionan los factores de crecimiento utilizados para calcular los exponentes de Lyapunov.
    }
    \label{fig:benettin_gram_schmidt}
\end{figure}
\shorthandon{>}

Después de \(N\) intervalos, el exponente de Lyapunov máximo viene dado por:
\begin{equation}
    \boxed{ \widehat{\lambda}_1(N)= \frac{1}{N\Delta T} \sum_{k=0}^{N-1}\ln g_k.}
    \label{eq:benettin-max}
\end{equation}

Así, el exponente es el resultado de la media aritmética del factor de crecimiento \textbf{renormalizado} \cite{benettin1980}. La renormalización \textit{mantiene la perturbación en el régimen lineal}. El valor de \(\Delta T\) no debe ser tan grande que la separación alcance escalas no lineales, ni tan pequeño que dominen errores de redondeo. La distancia inicial y el tiempo de renormalización deben someterse a pruebas de convergencia \cite{dubeibe2014}.

El Ejemplo~\ref{ej:benettin-intervalo} muestra el cálculo explícito correspondiente a un único intervalo de renormalización.

\begin{ejemplo}[Cálculo en un intervalo con $\Delta T = 0.01$]
\label{ej:benettin-intervalo}
Sea \(X_0=(1,1,1)^{\trans}\), entonces:
\begin{equation}
    J(X_0)=
    \begin{pmatrix}
        -10&10&0\\
        27&-1&-1\\
        1&1&-\frac83
    \end{pmatrix}.
\end{equation}
Con \(\Delta T=0.01\), una aproximación temporal de primer orden resulta en:
\begin{equation}
    \widetilde{\delta X}_1 \simeq \left[I+\Delta T\,J(X_0)\right]\delta X_0 = 10^{-6}
    \begin{pmatrix}
        0.90\\0.27\\0.01
    \end{pmatrix}.
\end{equation}
Luego,
\begin{equation}
    \norm{\widetilde{\delta X}_1} \simeq0.93968\times10^{-6}, \qquad g_0\simeq0.93968,
\end{equation}
y la tasa local de crecimiento es:
\begin{equation}
    \lambda_{\mathrm{loc},0} = \frac{1}{0.01}\ln(0.93968) \simeq-6.22.
\end{equation}
Obsérvese que \(\lambda_{\mathrm{loc},0}<0\): la tasa de un único
intervalo puede ser negativa; el exponente asintótico es el promedio
\eqref{eq:benettin-max} sobre \(N\) grande y resulta positivo.
Tras renormalizar:
\begin{equation}
    \delta X_1 \simeq10^{-6}
    \begin{pmatrix}
        0.9578\\0.2873\\0.0106
    \end{pmatrix}.
\end{equation}
Este procedimiento se repite durante \(N\) intervalos y el promedio
aritmético \eqref{eq:benettin-max} converge al exponente máximo de
Lyapunov global del sistema estudiado.
\end{ejemplo}
El diagrama de la Figura~\ref{fig:benettin_ciclo} resume el ciclo
completo del algoritmo para el exponente máximo.

\begin{figure}[htb]
    \centering
    \begin{tikzpicture}[
        node distance=1.1cm and 0.8cm,
        box/.style={draw,rounded corners,align=center,minimum height=1.15cm, minimum width=3.1cm,fill=blue!5},
        arr/.style={-{Latex[length=2.5mm]},thick}
    ]
        \node[box] (start) {$X(t_k),\,\delta X_k$\\norma \(\varepsilon\)};
        \node[box,right=of start] (prop) {Propagación en \(\Delta T\)\\ \(\dot X=F(X)\), \(\delta\dot X=J\delta X\)};
        \node[box,right=of prop] (growth) {Crecimiento\\ \(g_k=\norm{\widetilde{\delta X}_{k+1}}/\varepsilon\)};
        \node[box,below=of growth] (ren) {Renormalización\\ \(\delta X_{k+1}=\varepsilon \widetilde{\delta X}_{k+1}/ \norm{\widetilde{\delta X}_{k+1}}\)};
        \node[box,left=of ren] (sum) {Acumulación\\ \(\sum_k\ln g_k\)};
        \draw[arr] (start)--(prop);
        \draw[arr] (prop)--(growth);
        \draw[arr] (growth)--(ren);
        \draw[arr] (ren)--(sum);
        \draw[arr] (sum.west) -| (start.south);
    \end{tikzpicture}
    \caption{Diagrama de flujo del método de Benettin para el exponente máximo.}
    \label{fig:benettin_ciclo}
\end{figure}

\section{Espectro completo}
\label{sec:espectro-completo}

Para obtener los tres exponentes se propaga una matriz tangente cuyas
columnas son tres vectores de perturbación independientes,
\begin{equation}
    W(t)=
    \begin{pmatrix}
        |&|&|\\
        w_1(t)&w_2(t)&w_3(t)\\
        |&|&|
    \end{pmatrix},
    \qquad \dot W=J(X(t))W,\qquad W(0)=I.
\end{equation}
La condición inicial \(W(0)=I\) es genérica: cualquier base ortonormal
inicial sirve salvo un conjunto de medida nula, y la base identidad es
la elección más sencilla.

Cada intervalo propaga la base reortonormalizada del intervalo anterior.
Si \(Q_k\) es la base ortonormal al tiempo \(t_k\) (con \(Q_0=I\)),
la propagación durante \(\Delta T\) produce
\begin{equation}
    \widetilde W_{k+1} = \Phi(t_{k+1},t_k)\,Q_k,
    \label{eq:propaga-base}
\end{equation}
donde \(\Phi\) es el propagador de la ecuación variacional
\(\dot W=J(X(t))W\). Al final del intervalo se factoriza
\begin{equation}
    \widetilde W_{k+1}=Q_{k+1}R_{k+1},
\end{equation}
donde \(Q_{k+1}^{\trans}Q_{k+1}=I\) y \(R_{k+1}\) es triangular superior.
Se toma \(W_{k+1}=Q_{k+1}\) como base inicial del siguiente intervalo y
se acumulan los logaritmos de los elementos diagonales:
\begin{equation}
    \boxed{
    \widehat{\lambda}_i= \frac{1}{N\Delta T} \sum_{k=0}^{N-1}\ln\abs{(R_{k+1})_{ii}}, \qquad i=1,2,3.}
    \label{eq:benettin-espectro}
\end{equation}
Los elementos diagonales \((R_{k+1})_{ii}\) miden el estiramiento a lo
largo de cada dirección ortogonalizada sucesiva, de modo que el
estimador \eqref{eq:benettin-espectro} converge, para \(N\) grande, al
espectro ordenado
\begin{equation}
    \lambda_1 \ge \lambda_2 \ge \lambda_3,
\end{equation}

El primer exponente
\(\widehat{\lambda}_1\) coincide con el estimador
\eqref{eq:benettin-max} del capítulo anterior. Nótese que \((R_{k+1})_{ii}\) puede resultar negativo según el convenio
de signos de la factorización QR empleada; por ello se toma el valor
absoluto antes del logaritmo. El significado geométrico es el mismo
ilustrado en la Figura~\ref{fig:benettin_gram_schmidt}: la dinámica
tangente alinea los vectores con las direcciones dominantes sucesivas y
la reortonormalización restaura una base ortonormal, conservando solo la
información de estiramiento en la diagonal de \(R\).

La propagación tangente seguida de QR evita tanto el crecimiento o
colapso incontrolado de las normas como la pérdida de independencia
lineal. En aritmética finita, QR mediante reflexiones de Householder
suele conservar mejor la ortogonalidad que el Gram--Schmidt clásico \cite{gatica2026}




% =========================================================
\begin{thebibliography}{99}
% =========================================================

\bibitem{Ehrendorfer2003} M. Ehrendorfer
\emph{The Liouville Equation: Background - Historical Background}
The Liouville Equation in Atmospheric Predictability, 48–49. 2003.
\url{https://www.ecmwf.int/sites/default/files/elibrary/2003/9271-liouville-equation-atmospheric-predictability.pdf}

\bibitem{lorenz1963} E. N. Lorenz, \emph{Deterministic Nonperiodic Flow}, Journal of the Atmospheric Sciences, 20, 130--141, 1963. \url{https://doi.org/10.1175/1520-0469(1963)020%3C0130:DNF%3E2.0.CO;2}

\bibitem{benettin1980} G. Benettin, L. Galgani, A. Giorgilli y J.-M. Strelcyn, \emph{Lyapunov Characteristic Exponents for Smooth Dynamical Systems and for Hamiltonian Systems: A Method for Computing All of Them. Part 1: Theory}, Meccanica, 15, 9--20, 1980. \url{https://doi.org/10.1007/BF02128236}

\bibitem{dubeibe2014} F. L. Dubeibe y L. D. Bermúdez-Almanza, \emph{Optimal Conditions for the Numerical Calculation of the Largest Lyapunov Exponent for Systems of Ordinary Differential Equations}, International Journal of Modern Physics C, 25, 1450024, 2014. \url{https://arxiv.org/abs/1312.5970}

\bibitem{skokos2010} C. Skokos, \emph{The Lyapunov Characteristic Exponents and Their Computation}, Lecture Notes in Physics, 790, 63--135, 2010. \url{https://arxiv.org/abs/0811.0882}

\bibitem{semenov2022} M. E. Semenov, S. V. Borzunov y P. A. Meleshenko, \emph{A New Way to Compute the Lyapunov Characteristic Exponents for Non-Smooth and Discontinuous Dynamical Systems}, Nonlinear Dynamics, 2022. \url{https://doi.org/10.1007/s11071-022-07492-6}

\bibitem{gatica2026} T. Gatica, R. Haqshenas y E. van 't Wout, \emph{Robust Numerical Calculation of Lyapunov Exponents for Chaotic Bubble Dynamics}, Physics of Fluids, 38, 078130, 2026. \url{https://doi.org/10.1063/5.0334228}

\bibitem{kuznetsov2020} N. V. Kuznetsov, T. N. Mokaev, O. A. Kuznetsova y E. V. Kudryashova, \emph{The Lorenz System: Hidden Boundary of Practical Stability and the Lyapunov Dimension}, Nonlinear Dynamics, 102, 713--732, 2020. \url{https://doi.org/10.1007/s11071-020-05856-4}

\bibitem{sprott2005} J. C. Sprott, \emph{Lyapunov Exponent and Dimension of the Lorenz Attractor}, University of Wisconsin--Madison, revisión de 2005. \url{https://sprott.physics.wisc.edu/chaos/lorenzle.htm}

\bibitem{sysidentpy} W. Rocha, \emph{Lorenz System}, SysIdentPy Documentation. \url{https://sysidentpy.org/es/user-guide/tutorials/chaotic-systems/lorenz-system/}

\bibitem{larsson2023} E. Larsson y V. Ström, \emph{Numerical Solutions and Parameter Sensitivity of the Lorenz System}, KTH Royal Institute of Technology, 2023. \url{https://www.diva-portal.org/smash/get/diva2:1777317/FULLTEXT01.pdf}

\bibitem{zhou2021} S. Zhou y X. Wang, \emph{Simple estimation method for the largest Lyapunov exponent of continuous fractional-order differential equations}, Physica A: Statistical Mechanics and its Applications, 563, 125478, 2021. \url{https://doi.org/10.1016/j.physa.2020.125478}

\bibitem{yakovleva2025} T. V. Yakovleva, A. V. Krysko, V. V. Dobriyan y V. A. Krysko, \emph{A modified neural network method for computing the Lyapunov exponent spectrum in the nonlinear analysis of dynamical systems}, Communications in Nonlinear Science and Numerical Simulation, 140, 108397, 2025. \url{https://doi.org/10.1016/j.cnsns.2024.108397}

\end{thebibliography}

\appendix{}

\chapter{Gram--Schmidt dentro del método de Benettin}
\label{app:gram-schmidt}
% =========================================================

Sean \(\widetilde w_1,\widetilde w_2,\widetilde w_3\) los vectores tangentes tras un intervalo. Gram--Schmidt construye una base ortonormal \(q_1,q_2,q_3\) por proyecciones sucesivas. Para el primer vector basta normalizar:
\begin{align}
    u_1&=\widetilde w_1, & r_{11}&=\norm{u_1}, & q_1&=\frac{u_1}{r_{11}}.
\end{align}
El segundo vector se limpia de su componente a lo largo de \(q_1\) y se normaliza:
\begin{align}
    r_{12}&=\langle q_1,\widetilde w_2\rangle, & u_2&=\widetilde w_2-r_{12}q_1,\\
    r_{22}&=\norm{u_2}, & q_2&=\frac{u_2}{r_{22}}.
\end{align}
El tercero se limpia de sus componentes a lo largo de \(q_1\) y \(q_2\):
\begin{align}
    r_{13}&=\langle q_1,\widetilde w_3\rangle, &
    r_{23}&=\langle q_2,\widetilde w_3\rangle,\\
    u_3&=\widetilde w_3-r_{13}q_1-r_{23}q_2, &
    r_{33}&=\norm{u_3}, &
    q_3&=\frac{u_3}{r_{33}}.
\end{align}
En general,
\begin{equation}
    u_j=\widetilde w_j- \sum_{i=1}^{j-1} \langle q_i,\widetilde w_j\rangle q_i, \qquad q_j=\frac{u_j}{\norm{u_j}}.
\end{equation}
El procedimiento está bien definido siempre que los \(\widetilde w_j\)
sean linealmente independientes (\(r_{jj}>0\)), lo cual se cumple
genéricamente si \(\Delta T\) no es excesivo.

Nótese que, por construcción, \(r_{jj}=\norm{u_j}>0\): la diagonal de
\(R\) es positiva en este convenio de Gram--Schmidt. Otros convenios de
factorización QR pueden producir diagonales negativas, de ahí el valor
absoluto en el estimador \eqref{eq:benettin-espectro}.

Sin ortogonalización, todos los vectores tienden a alinearse con la dirección de \(\lambda_1\). Las proyecciones sucesivas eliminan las componentes ya atribuidas a los subespacios de crecimiento superior. Los productos
\[
    r_{11},\qquad r_{11}r_{22},\qquad r_{11}r_{22}r_{33}
\]
miden, respectivamente, crecimiento de longitud, área y volumen tangentes: como \(R\) es triangular superior, el determinante de cada paso es el producto de su diagonal, \(\det\widetilde W_k=\prod_i r_{ii}^{(k)}\), y el determinante del propagador compuesto es el producto de los determinantes de cada intervalo. Por ello,
\begin{equation}
    \lambda_1+\lambda_2+\lambda_3 = \lim_{T\to\infty}\frac1T \ln\abs{\det\Phi(T,0)},
    \label{eq:suma}
\end{equation}
lo que recupera la identidad de la suma de exponentes. Para el sistema de Lorenz, \(\operatorname{tr}J=-(\sigma+1+\beta)=-41/3\), de modo que \eqref{eq:suma} proporciona un control cerrado del cálculo numérico: la suma de los tres exponentes debe aproximarse a \(-13.67\).


\chapter{Obtención algebraica de los puntos de equilibrio}
\label{app:equilibrios}
% =========================================================

Los equilibrios satisfacen \(F(X_e)=0\):
\begin{equation}
    \sigma(y_e-x_e)=0,\qquad x_e(\rho-z_e)-y_e=0,\qquad x_ey_e-\beta z_e=0.
\end{equation}

Como \(\sigma\neq0\),
\begin{equation}
    y_e=x_e.
\end{equation}

La segunda ecuación queda
\begin{equation}
    x_e(\rho-z_e-1)=0.
\end{equation}

Si \(x_e=0\), entonces \(y_e=0\) y la tercera ecuación implica \(z_e=0\); por tanto,
\begin{equation}
    E_0=(0,0,0).
\end{equation}

Si \(x_e\neq0\), se tiene
\begin{equation}
    z_e=\rho-1.
\end{equation}

La tercera ecuación proporciona
\begin{equation}
    x_e^2=\beta(\rho-1),
\end{equation}
y, para \(\rho>1\),
\begin{equation}
    E_\pm= \left( \pm\sqrt{\beta(\rho-1)}, \pm\sqrt{\beta(\rho-1)}, \rho-1 \right).
\end{equation}

Este procedimiento determina la localización geométrica de los estados estacionarios. Su clasificación dinámica requiere posteriormente resolver el problema espectral
\[
    \det\!\left(J(X_e)-\lambda I\right)=0.
\]


\end{document}